% !TeX program = xelatex
\documentclass[aspectratio=169,11pt]{beamer}
\usepackage{fontspec}
\setsansfont{Liberation Sans}
\setmainfont{Liberation Sans}
\setmonofont{Liberation Mono}
\usepackage{polyglossia}
\setdefaultlanguage{russian}
\setotherlanguage{english}
\usepackage{amsmath,booktabs,tabularx,array,graphicx,tikz,pgfplots}
\usetikzlibrary{arrows.meta,positioning,calc}
\pgfplotsset{compat=1.18}
\usefonttheme{professionalfonts}
\setbeamertemplate{navigation symbols}{}
\setbeamertemplate{headline}{}
\setbeamertemplate{frametitle}{\vspace{2mm}\centering\bfseries\insertframetitle\par}
\setbeamercolor{normal text}{fg=black,bg=white}
\setbeamercolor{frametitle}{fg=black,bg=white}
\setbeamercolor{itemize item}{fg=black}
\setbeamercolor{enumerate item}{fg=black}
\setbeamertemplate{itemize item}{\raisebox{0.2ex}{\scriptsize$\bullet$}}
\setbeamertemplate{itemize subitem}{--}
\setbeamertemplate{footline}{\hfill\usebeamerfont{page number in head/foot}\insertframenumber/\inserttotalframenumber\hspace{5mm}\vspace{3mm}}
\setbeamersize{text margin left=7mm,text margin right=7mm}
\setlength{\parskip}{3pt}
\setbeamertemplate{note page}[plain]
\newcommand{\src}[1]{\par\vspace{1mm}{\fontsize{7}{8}\selectfont Источник: Герасименко А.\,Ю. и др. Известия вузов. Электроника. 2010. №~4(84). #1\par}}
\newcommand{\smalltext}{\fontsize{9}{10.8}\selectfont}
\newcommand{\article}{8102.pdf}
\newcommand{\sourcepanel}[3]{%
\IfFileExists{\article}{\includegraphics[page=#1,width=\linewidth,height=#2,keepaspectratio]{\article}}{%
\fbox{\parbox[c][#2][c]{0.93\linewidth}{\centering\small #3\\[2mm]\footnotesize Для воспроизведения рисунка поместите\\\texttt{8102.pdf} рядом с исходником.}}}}
\tikzset{box/.style={draw,rectangle,align=center,inner sep=4pt,font=\small},arr/.style={-{Stealth[length=2mm]},thick}}
\begin{document}

% Титульный — первый слайд, учитывается в общем числе, без номера.
\begin{frame}[plain]
\begin{tikzpicture}[remember picture,overlay]
\node[anchor=north,align=center,text width=15.6cm,font=\fontsize{10.5}{12}\selectfont]
 at ([yshift=-3mm]current page.north) {МИНОБРНАУКИ РОССИИ\\[2pt]
Федеральное государственное автономное образовательное учреждение высшего образования\\
«Национальный исследовательский университет\\
«Московский институт электронной техники»};
\node[align=center,font=\fontsize{10.5}{12}\selectfont] at ([yshift=-20.8mm]current page.north) {Институт биомедицинских систем};
\node[align=center,font=\fontsize{17}{20}\selectfont] at ([yshift=-33.8mm]current page.north) {Гуслянцев Егор Евгеньевич};
\node[align=center,text width=14.8cm,font=\bfseries\fontsize{17}{20}\selectfont] at ([yshift=-52mm]current page.north) {Нанокомпозитный припой для лазерной сварки\\биологических тканей};
\node[anchor=east,font=\fontsize{10.5}{12}\selectfont] at ([xshift=-9mm,yshift=13.8mm]current page.south east) {Руководитель: к.ф.-м.н., доцент Потапов Д.А.};
\node[font=\fontsize{10.5}{12}\selectfont] at ([yshift=4mm]current page.south) {Москва 2026};
\end{tikzpicture}
\end{frame}

\begin{frame}{Задача исследования и принцип соединения}
{\centering\small Нанокомпозитный припой для лазерной сварки биологических тканей\par}
\vspace{1mm}
\begin{itemize}
\item Задача: сравнить прочность соединений при изменении углеродного наполнителя в БСА.
\item Объекты: свиная кожа и хрящ бычьей трахеи; \textit{in vitro}.
\end{itemize}
\vspace{2mm}
\begin{center}
\begin{tikzpicture}[node distance=3mm]
\node[box,text width=29mm](a){Нанесение\\припоя};
\node[box,text width=29mm,right=of a](b){Сопоставление\\краёв};
\node[box,text width=29mm,right=of b](c){Облучение\\970 нм};
\node[box,text width=25mm,right=of c](d){Соединение};
\draw[arr](a)--(b);\draw[arr](b)--(c);\draw[arr](c)--(d);
\begin{scope}[shift={(5,-1.5)}]
\draw[fill=black!8](-2,0)rectangle(-0.08,0.45);
\draw[fill=black!8](0.08,0)rectangle(2,0.45);
\draw[fill=black!45](-0.08,0)rectangle(0.08,0.65);
\draw[arr](0,1.3)--(0,0.7);
\node[anchor=west,font=\scriptsize] at (0.2,1.1){Лазерное излучение};
\node[anchor=east,font=\scriptsize] at (-2.1,0.2){Ткань};
\node[anchor=west,font=\scriptsize] at (2.1,0.2){Ткань};
\node[font=\scriptsize] at (0,-0.22){Слой припоя};
\end{scope}
\end{tikzpicture}
\end{center}
\small БСА (бычий сывороточный альбумин) — белковая матрица.\\
УНТ (углеродные нанотрубки) — поглощение излучения и предполагаемое упрочнение.\\
Температурный диапазон, указанный авторами как оптимальный: 65–70~°C.
\par\vspace{1mm}{\fontsize{6.5}{7.5}\selectfont
Источник: Герасименко А.\,Ю., Губарьков О.\,В., Ичкитидзе Л.\,П., Подгаецкий В.\,М., Селищев С.\,В., Пономарева О.\,В.\\
Нанокомпозитный припой для лазерной сварки биологических тканей // Известия вузов. Электроника. 2010. №~4(84). С.~33–41.\par}
\note{Схема условная. Нанотрубочный каркас не изображается как наблюдённая структура. Безрубцовое заживление, стерильность и клиническое превосходство в работе не проверены.}
\end{frame}

\begin{frame}{Составы и приготовление припоев}
\centering Основная прочностная серия: 28 мас.\,\% БСА;\\0,1 мас.\,\% наполнителя при его наличии.
\vspace{2mm}
\smalltext
\begin{tabularx}{\linewidth}{@{}p{39mm}p{31mm}X@{}}
\toprule Наполнитель & Обозначение & Характеристика\\\midrule
Без наполнителя & БСА & Альбуминовый контроль\\
Сажа К-354 & БСА + сажа & Углеродный контроль\\
Активированный уголь & БСА + МРАУ & Средний размер частиц $\leq 10$ мкм\\
Многослойные УНТ МИЭТ & БСА + МУНТ I & Наружный диаметр 3–30 нм\\
Многослойные УНТ «Таунит» & БСА + МУНТ II & Диаметр 20–40 нм; средняя длина $\geq 2$ мкм\\
Однослойные УНТ & БСА + ОУНТ & Средний диаметр $\sim 1$ нм; длина $\sim 1$ мкм\\\bottomrule
\end{tabularx}
\vspace{3mm}
\begin{tikzpicture}[node distance=2mm]
\node[box,text width=19mm,font=\scriptsize](a){БСА + вода};
\node[box,text width=23mm,font=\scriptsize,right=of a](b){Механическое перемешивание};
\node[box,text width=35mm,font=\scriptsize,right=of b](c){Декантирование, выдержка, фильтрование};
\node[box,text width=21mm,font=\scriptsize,right=of c](d){Введение наполнителя};
\node[box,text width=26mm,font=\scriptsize,right=of d](e){УЗ-обработка; контроль температуры};
\foreach \x/\y in {a/b,b/c,c/d,d/e}\draw[arr](\x)--(\y);
\end{tikzpicture}
\par\vspace{2mm}\raggedright\scriptsize
Описанная процедура приготовления относится к 25\%-ной дисперсии БСА,\newline а не к составу основной прочностной серии.
\src{С.~35–36, 39.}
\note{25\%-ная дисперсия: 8,33 г БСА на 25 мл воды; 3 ч перемешивания при температуре не выше 30 °C; выдержка 12 ч. Способ перехода к 28\%-ному составу не реконструирован. Спектральные иллюстрации с 0,3 мас.\,\% УНТ относятся к другой серии.}
\end{frame}

\begin{frame}{Оптические свойства и стабильность}
\begin{columns}[T,onlytextwidth]
\begin{column}{0.49\linewidth}
\small Спектры поглощения припоев с УНТ
\par\vspace{1mm}
% Полная исходная страница сохраняет оси, подписи и номера кривых.
% Для показа крупным планом можно заменить на отдельный файл fig4.pdf.
\IfFileExists{fig4.pdf}{\includegraphics[width=\linewidth,height=43mm,keepaspectratio]{fig4.pdf}}{\sourcepanel{5}{43mm}{Рис. 4 статьи: спектры поглощения}}
\par\scriptsize УНТ: 0,3 мас.\,\%; слой 1 мм.\\
Кривые: 1 — «НаноКарбонЛайт» (ОУНТ);\\2 — «Таунит» (МУНТ II); 3 — МИЭТ (МУНТ I).
\end{column}
\begin{column}{0.49\linewidth}
\smalltext
\begin{tabularx}{\linewidth}{@{}p{18mm}X@{}}
\toprule Испытание & Условия и наблюдения\\\midrule
Термо- & 936 ч при комнатной температуре; отдельное испытание при 42~°C. Спектры изменились незначительно.\\[2mm]
Фото- & Светодиод $\sim$630 нм; 65 мВт; 83 ч. Небольшие изменения спектров после облучения.\\\bottomrule
\end{tabularx}
\vspace{2mm}
\scriptsize Серия испытаний стабильности: 25\% БСА + 0,3 мас.\,\% «Таунит».\\[2mm]
При 42~°C: 85 ч в тексте; 83 ч в подписи к рис.~1.\\[2mm]
\small Оценивались изменения спектров, а не биологическая безопасность.
\end{column}
\end{columns}
\src{С.~36–38, рис.~1, 3, 4.}
\note{Поглощение ОУНТ ниже, чем «Таунит», а «Таунит» ниже, чем МИЭТ; авторы связывают различия со степенью очистки. Это не доказательство механизма упрочнения. Для резервного обсуждения использовать рис. 1 и 3: слой на рис. 1 — 1 мм; на рис. 3а — 0,1 мм, на рис. 3б — 1 мм. Светодиод 630 нм не эквивалентен сварочному лазеру 970 нм.}
\end{frame}

\begin{frame}{Лазерная сварка и измерение прочности}
\begin{columns}[T,onlytextwidth]
\begin{column}{0.47\linewidth}
\small Облучение хряща и вид соединения
\par\vspace{1mm}
\IfFileExists{fig5.pdf}{\includegraphics[width=\linewidth,height=36mm,keepaspectratio]{fig5.pdf}}{\sourcepanel{6}{36mm}{Рис. 5 статьи: облучение и лазерный шов}}
\par\scriptsize Иллюстративная серия: 25\% БСА +\\0,3 мас.\,\% «Таунит». Чёрный цвет обусловлен припоем.
\end{column}
\begin{column}{0.51\linewidth}
\smalltext Образцы: 30–33 $\times$ 2–4 $\times$ 0,5–1,5 мм; хранение в формалине.\\[1mm]
Разрез $\rightarrow$ припой $\rightarrow$ фиксация $\rightarrow$ облучение.\\[1mm]
Непрерывный диодный лазер: 970 нм;\\5–15 Вт на выходе волокна; 2–5 мин.\\[1mm]
Контроль ИК-пирометром: 65–70~°C.
\begin{center}
\begin{tikzpicture}[x=0.82cm,y=0.8cm]
\draw[fill=black!30](-2,0)rectangle(-1.6,0.8);
\draw[fill=black!30](1.6,0)rectangle(2,0.8);
\draw[fill=black!8](-1.6,0.2)rectangle(1.6,0.6);
\draw[very thick](0,0.2)--(0,0.6);
\draw[arr](-2,0.4)--(-2.9,0.4)node[left]{$F$};
\draw[arr](2,0.4)--(2.9,0.4)node[right]{$F$};
\draw[arr](0.7,1.2)node[above,font=\scriptsize]{Сечение $S$}--(0,0.65);
\node[font=\scriptsize] at(0,-0.2){Растяжение образца};
\end{tikzpicture}
\end{center}
\vspace{-2mm}
\[\sigma=\frac{F_{\max}}{S}\]
\vspace{-3mm}\scriptsize Не менее пяти образцов на вариант.\\Контроль — неразрезанная ткань.
\end{column}
\end{columns}
\src{С.~37–39.}
\note{Формула даёт номинальное напряжение при разрушении, а не силу. Плотность мощности без площади пятна не вычисляется. В тексте фотография названа рис. 5; в подписи исходного PDF напечатано «Рис. 4». Здесь сохранена нумерация slides.md, это редакционная несогласованность источника.}
\end{frame}

\begin{frame}{Влияние наполнителя на прочность}
\centering\small 28 мас.\,\% БСА; 0,1 мас.\,\% наполнителя при его наличии.
\vspace{1mm}
\begin{columns}[T,onlytextwidth]
\begin{column}{0.5\linewidth}
\centering\small Свиная кожа
\begin{tikzpicture}
\begin{axis}[width=7.45cm,height=4.4cm,xmin=0,xmax=2.5,ymin=-0.6,ymax=5.6,
 ytick={0,1,2,3,4,5},yticklabels={БСА,БСА + сажа,БСА + МРАУ,БСА + МУНТ I,БСА + МУНТ II,БСА + ОУНТ},y dir=reverse,
 tick label style={font=\tiny},xlabel={\scriptsize Прочность, МПа},xtick={0,0.5,1,1.5,2,2.5},
 axis x line*=bottom,axis y line*=left,xmajorgrids,grid style={black!12},bar width=9pt,bar shift=0pt,
 point meta=explicit symbolic,nodes near coords, nodes near coords align={horizontal},
 every node near coord/.append style={font=\tiny,anchor=west,xshift=12pt},
 error bars/x dir=both,error bars/x explicit]
\addplot[xbar,fill=black!25,draw=black] coordinates {(0.4,0) +- (0.1,0) [0,4]};
\addplot[xbar,fill=black!45,draw=black] coordinates {(0.3,1) +- (0.1,0) [0,3] (0.2,2) +- (0.1,0) [0,2]};
\addplot[xbar,fill=black!75,draw=black] coordinates {(1.4,3) +- (0.3,0) [1,4] (1.2,4) +- (0.3,0) [1,2] (0.8,5) +- (0.3,0) [0,8]};
\end{axis}
\end{tikzpicture}
\end{column}
\begin{column}{0.5\linewidth}
\centering\small Бычий хрящ
\begin{tikzpicture}
\begin{axis}[width=7.45cm,height=4.4cm,xmin=0,xmax=2.5,ymin=-0.6,ymax=5.6,
 ytick={0,1,2,3,4,5},yticklabels={БСА,БСА + сажа,БСА + МРАУ,БСА + МУНТ I,БСА + МУНТ II,БСА + ОУНТ},y dir=reverse,
 tick label style={font=\tiny},xlabel={\scriptsize Прочность, МПа},xtick={0,0.5,1,1.5,2,2.5},
 axis x line*=bottom,axis y line*=left,xmajorgrids,grid style={black!12},bar width=9pt,bar shift=0pt,
 point meta=explicit symbolic,nodes near coords,nodes near coords align={horizontal},
 every node near coord/.append style={font=\tiny,anchor=west,xshift=14pt},
 error bars/x dir=both,error bars/x explicit]
\addplot[xbar,fill=black!25,draw=black] coordinates {(0.9,0) +- (0.3,0) [0,9]};
\addplot[xbar,fill=black!45,draw=black] coordinates {(0.6,1) +- (0.2,0) [0,6] (0.5,2) +- (0.2,0) [0,5]};
\addplot[xbar,fill=black!75,draw=black] coordinates {(1.9,3) +- (0.5,0) [1,9] (1.4,4) +- (0.4,0) [1,4] (1.7,5) +- (0.5,0) [1,7]};
\end{axis}
\end{tikzpicture}
\end{column}
\end{columns}
\vspace{-2mm}\raggedright\small УНТ повышают среднюю прочность относительно БСА; сажа и МРАУ её снижают.
\par\scriptsize Интервалы $\pm$ воспроизведены из статьи; их статистический смысл не указан.\\Статистическое превосходство отдельных нанотрубочных составов не установлено.
\src{С.~39, таблица.}
\note{Не обозначать интервалы как стандартное отклонение, стандартную ошибку или доверительный интервал. Проверка значимости в статье не представлена. Данные цитируемых работ с другими тканями и лазерами не включены.}
\end{frame}

\begin{frame}{Сравнение с прочностью исходной ткани}
\centering\small Доля прочности контрольной ткани: $100\,\sigma/\sigma_{\text{т}}$.
\vspace{1mm}
\begin{columns}[T,onlytextwidth]
\begin{column}{0.5\linewidth}
\centering\small Свиная кожа
\begin{tikzpicture}
\begin{axis}[width=7.3cm,height=4.3cm,xmin=0,xmax=100,ymin=-0.6,ymax=2.6,y dir=reverse,
 ytick={0,1,2},yticklabels={БСА,БСА + МУНТ I,Ткань},tick label style={font=\scriptsize},
 xlabel={\scriptsize Доля прочности ткани, \%},xtick={0,25,50,75,100},axis x line*=bottom,axis y line*=left,
 xmajorgrids,grid style={black!12},bar width=15pt,bar shift=0pt]
\addplot[xbar,fill=black!25] coordinates {(2.7,0)};
\addplot[xbar,fill=black!75] coordinates {(9.3,1)};
\addplot[xbar,fill=white,draw=black] coordinates {(100,2)};
\node[anchor=west,font=\tiny] at(axis cs:5,0){2,7\%; $0{,}4\pm0{,}1$ МПа};
\node[anchor=west,font=\tiny] at(axis cs:12,1){9,3\%; $1{,}4\pm0{,}3$ МПа};
\node[font=\tiny] at(axis cs:53,2){100\%; $15\pm5$ МПа};
\end{axis}
\end{tikzpicture}
\end{column}
\begin{column}{0.5\linewidth}
\centering\small Бычий хрящ
\begin{tikzpicture}
\begin{axis}[width=7.3cm,height=4.3cm,xmin=0,xmax=100,ymin=-0.6,ymax=2.6,y dir=reverse,
 ytick={0,1,2},yticklabels={БСА,БСА + МУНТ I,Ткань},tick label style={font=\scriptsize},
 xlabel={\scriptsize Доля прочности ткани, \%},xtick={0,25,50,75,100},axis x line*=bottom,axis y line*=left,
 xmajorgrids,grid style={black!12},bar width=15pt,bar shift=0pt]
\addplot[xbar,fill=black!25] coordinates {(15,0)};
\addplot[xbar,fill=black!75] coordinates {(32,1)};
\addplot[xbar,fill=white,draw=black] coordinates {(100,2)};
\node[anchor=west,font=\tiny] at(axis cs:18,0){15\%; $0{,}9\pm0{,}3$ МПа};
\node[anchor=west,font=\tiny] at(axis cs:35,1){32\%; $1{,}9\pm0{,}5$ МПа};
\node[font=\tiny] at(axis cs:53,2){100\%; $6\pm2$ МПа};
\end{axis}
\end{tikzpicture}
\end{column}
\end{columns}
\vspace{1mm}\raggedright\small Упрочнение относительно БСА не означает достижения прочности исходной ткани.
\par\scriptsize Проценты приведены по таблице статьи. Контроль — неразрезанная ткань,\\а не соединение хирургическим швом. Заживление не исследовалось.
\src{С.~39, таблица.}
\end{frame}

\begin{frame}{Механизм: наблюдение и гипотеза}
\begin{columns}[T,onlytextwidth]
\begin{column}{0.48\linewidth}
\centering\textbf{Экспериментальные наблюдения}\par\raggedright
\begin{itemize}
\item Более высокие средние значения прочности при введении УНТ.
\item Отсутствие аналогичного эффекта у сажи и МРАУ.
\item Различный порядок нанотрубочных составов для двух тканей.
\end{itemize}
\end{column}
\begin{column}{0.49\linewidth}
\centering\textbf{Гипотеза авторов}\par
\begin{tikzpicture}[node distance=3mm]
\node[box,text width=58mm,font=\small](a){Лазерная обработка};
\node[box,text width=58mm,font=\small,below=of a](b){Структурирование УНТ};
\node[box,text width=58mm,font=\small,below=of b](c){Каркас с взаимодействующим альбумином};
\node[box,text width=58mm,font=\small,below=of c](d){Упрочнение слоя};
\foreach \x/\y in {a/b,b/c,c/d}\draw[arr,dashed](\x)--(\y);
\end{tikzpicture}
\end{column}
\end{columns}
\vspace{3mm}\small Прямого подтверждения каркаса в данной работе не представлено.
\par\small Механические измерения не разделяют вклад поглощения, нагрева и армирования.
\src{С.~39–40.}
\note{Схема является моделью, а не микрофотографией или результатом наблюдения самосборки. Авторы отмечают возможную недостаточность предложенной физической картины.}
\end{frame}

\begin{frame}{Ограничения и результаты}
\smalltext
\begin{tabularx}{\linewidth}{@{}p{74mm}X@{}}
\toprule Ограничение & Следствие\\\midrule
\textit{In vitro}; ткани хранились в формалине & Нет прямого переноса на живые ткани\\[1mm]
Режим облучения менялся с составом & Сравнивается сочетание состава и обработки\\[1mm]
Не менее пяти образцов; смысл $\pm$ не уточнён; нет проверки значимости & Не установлено статистическое превосходство составов\\[1mm]
Нет оценки биологической реакции, герметичности, циклической и длительной прочности & Клиническая применимость не установлена\\\bottomrule
\end{tabularx}
\vspace{2mm}
\begin{enumerate}
\item УНТ повышают среднюю прочность относительно БСА в исследованных условиях.
\item Наибольшие средние значения — с МУНТ I: 1,4 МПа для кожи и 1,9 МПа для хряща; 9,3\% и 32\% прочности контрольных тканей.
\item Следующие задачи: сопоставимые режимы облучения, проверка структуры, исследования на живых тканях.
\end{enumerate}
\src{С.~37–40.}
\note{Отсутствие испытаний не доказывает небезопасность. Не утверждать клиническую готовность или превосходство над хирургическим швом. Слайд остаётся на экране во время вопросов.}
\end{frame}
\end{document}
